% !TeX root = ../tfg.tex
% !TeX encoding = utf8

\chapter{Marco experimental}
\label{ch:experimental}

Los experimentos de este trabajo se apoyan sobre dos ejes: el conjunto de datos sobre el que se evalúa la detección OOD y el tipo de modelo generativo empleado para la misma. Para ello, se entrena un modelo exclusivamente con muestras de la distribución ID, mientras que los datos OOD solo se utilizan en la fase de evaluación, como hemos acostumbrado en el resto de la memoria.

\section{Descripción de los datos}
\label{sec:exp-dataset}

Los conjuntos de datos empleados cubren un amplio espectro de dificultad. Estos van desde espacios sintéticos bidimensinoales con distinto nivel de separabilidad geométrica hasta espacios de características de alta dimensión extraídas con un modelo fundacional.

\paragraph{Blobs y Moons.}~\\
Los conjuntos Moons y Blobs se generan mediante las utilidades de \texttt{sklearn.datasets}~\cite{pedregosa_scikit-learn_2011}. En \textit{Blobs}, $N{=}2\,000$ puntos se generan desde
dos centroides en $(-2, 0)$ y $(2, 0)$ con desviación~$\sigma{=}0.5$; el
clúster~0 es la distribución ID y el clúster~1 la OOD. En \emph{Moons} se generan $N{=}2\,000$ puntos bidimensionales con
ruido $\sigma{=}0.05$: la clase~0, cuya geometría es la de una media luna, actúa
como distribución ID y la clase~1, adyacente y con forma cóncava complementaria,
como distribución OOD. En concreto, la separabilidad lineal de \emph{Blobs} contrasta con la estructura no lineal de \emph{Moons}, lo que permite evaluar si los modelos son sensibles tanto a geometrías regulares como a distribuciones con mayor complejidad topológica. Las
figuras~\ref{fig:input_space_moons} y~\ref{fig:input_space_blobs} muestran el espacio de entrada de ambos conjuntos.

\begin{figure}[htbp]
    \centering
    \includegraphics[width=\textwidth]{results-training-mlp_vae_moons_s42_lr0.001-plots-input_space_vs_ood.pdf}
    \caption{Proyección PCA del espacio de entrada para el conjunto sintético
             \textit{Moons}: clase~0 (ID, media luna superior) frente a clase~1
             (OOD, media luna inferior). La frontera no lineal entre ambas clases
             contrasta con la separabilidad lineal de \textit{Blobs}.}
    \label{fig:input_space_moons}
\end{figure}

\begin{figure}[htbp]
    \centering
    \includegraphics[width=\textwidth]{results-training-mlp_vae_blobs_s42_lr0.001-plots-input_space_vs_ood.pdf}
    \caption{Proyección PCA del espacio de entrada para el conjunto sintético
             \textit{Blobs}: clúster~0 en~$(-2,0)$ (ID) frente a clúster~1
             en~$(2,0)$ (OOD). La separabilidad lineal facilita la validación
             del comportamiento esperado de los modelos.}
    \label{fig:input_space_blobs}
\end{figure}

\paragraph{MNIST frente a SVHN.}~\\
MNIST~\cite{lecun_gradient-based_1998} comprende $60\,000$ imágenes de entrenamiento
y $10\,000$ de test de dígitos manuscritos en escala de grises ($28\times28$ píxeles)
y se emplea en su totalidad como distribución ID. Como distribución OOD se utiliza
SVHN~\cite{netzer_reading_2011}, un conjunto de dígitos fotorrealistas tomados de
fachadas de edificios en color, que se convierte a escala de grises y se redimensiona
a $28\times28$ para homogeneizar la dimensión de entrada con MNIST. Aunque ambos
conjuntos pertenecen al dominio semántico de los dígitos, la brecha entre el dominio
dibujado y normalizado de MNIST y el dominio fotorrealista con fondo natural de SVHN
genera un contraste entre dominios de dificultad moderada, ampliamente utilizado como
referencia en la literatura de detección OOD. En la
proyección PCA bidimensional del espacio de píxeles
(figura~\ref{fig:input_space_mnist}), ambas distribuciones presentan separación
apreciable, si bien con cierto solapamiento en la periferia de ambas nubes.

\begin{figure}[htbp]
    \centering
    \includegraphics[width=0.8\textwidth]{results-training-mlp_vae_mnist_s42_lr0.0001-plots-input_space_vs_ood.pdf}
    \caption{Proyección PCA del espacio de entrada para MNIST (ID) frente a
             SVHN~(OOD). Al provenir de dominios distintos, la proyección
             bidimensional muestra la separabilidad entre distribuciones en el
             espacio de píxeles.}
    \label{fig:input_space_mnist}
\end{figure}

\paragraph{PathMNIST.}~\\
PathMNIST~\cite{yang_medmnist_2023} es un subconjunto del \textit{benchmark}
MedMNIST v2 formado por parches histológicos de tejido colorrectal de $28\times28$
píxeles en formato RGB, clasificados en nueve clases de tejido. Se definen dos
configuraciones de partición ID/OOD de dificultad creciente:
\begin{itemize}
    \item \textbf{PathMNIST-C1}: distribución ID~$=$~clase~6 (\textit{normal colon
          mucosa}); OOD~$=$~clases~7 (\textit{stroma}) y~8 (\textit{colorectal
          adenocarcinoma epithelium}).
    \item \textbf{PathMNIST-C2}: distribución ID~$=$~clases~0, 3, 4, 5 y~6 (tejidos
          no tumorales, excluyendo las clases~1 (\textit{background}) y~2 (\textit{debris})); OOD~$=$~clases~7 y~8.
\end{itemize}
La elección de las clases~7 y~8 como OOD no es arbitraria: el estroma y el epitelio
tumoral comparten características morfológicas parciales con el tejido sano, lo que
convierte este escenario en un caso representativo de alta dificultad en detección OOD
de dominio médico. El solapamiento
entre distribuciones en el espacio de píxeles para la configuración C1 es visible en
la figura~\ref{fig:input_space_pathmnist}. La configuración C2 amplía la cobertura ID a cinco clases, lo que
aumenta la variabilidad interna de la distribución de entrenamiento. 

\begin{figure}[htbp]
    \centering
    \includegraphics[width=\textwidth]{results-training-mlp_vae_pathmnist_c1_s42_lr0.0001-plots-input_space_vs_ood.pdf}
    \caption{Proyección PCA del espacio de entrada para PathMNIST: tejido normal
             (ID, clase~6) frente a tejido tumoral (OOD, clase~8). El solapamiento
             entre distribuciones histológicas en el espacio de píxeles ilustra la
             dificultad del problema.}
    \label{fig:input_space_pathmnist}
\end{figure}

\paragraph{SICAP.}~\\
SICAP~\cite{silva-rodriguez_going_2020} es un conjunto de biopsias de próstata anotadas
con el sistema Gleason, que cuantifica la agresividad del cáncer de próstata en cuatro
grados (1--4). A diferencia del resto de conjuntos, las imágenes no se procesan
directamente como píxeles; en su lugar se emplean vectores de características de
dimensión~$1\,024$ extraídos con \textbf{UNI}~\cite{chen_towards_2024}, un \textit{vision
transformer} ViT-L/16 pre-entrenado mediante DINOv2 sobre más de cien millones de
parches histopatológicos procedentes de The Cancer Genome Atlas (TCGA) y otras
colecciones públicas. El uso de estas características sitúa a SICAP en el dominio de
alta dimensión estructurada, capturando información morfológica de alto nivel que los
píxeles crudos no podrían proporcionar de forma tratable computacionalmente. Se definen
dos configuraciones:
\begin{itemize}
    \item \textbf{SICAP-C1}: ID~$=$~Gleason~1 (tejido benigno estricto); OOD~$=$~Gleason~2,
          3 y~4.
    \item \textbf{SICAP-C12}: ID~$=$~Gleason~1 y~2 (bajo grado); OOD~$=$~Gleason~3 y~4
          (alto grado).
\end{itemize}
SICAP-C1 modela el escenario más restrictivo, en el que solo se dispone de tejido
completamente benigno durante el entrenamiento, mientras que SICAP-C12 amplía la
definición de normalidad al incluir cambios de bajo grado, reduciendo así la
distancia estadística entre ID y OOD. El solapamiento en el espacio de características
UNI entre grados adyacentes, visible en la figura~\ref{fig:input_space_sicap}, anticipa
la dificultad del problema respecto al resto de conjuntos.

\begin{figure}[htbp]
    \centering
    \includegraphics[width=0.7\textwidth]{results-training-mlp_vae_sicap_c12_s42_lr0.0001-plots-input_space_vs_ood.pdf}
    \caption{Proyección PCA de las características UNI del conjunto SICAP-C12:
             Gleason~1--2 (ID) frente a Gleason~3--4 (OOD). El solapamiento en
             el espacio de características refleja la dificultad de separar grados
             de Gleason mediante distancias simples.}
    \label{fig:input_space_sicap}
\end{figure}


\section{Procesamiento de los datos}
\label{sec:exp-pipeline}

\subsection{Partición de los conjuntos}

En todos los experimentos se respeta el principio fundamental de la evaluación OOD:
\textbf{el conjunto de entrenamiento contiene únicamente muestras ID}, y los datos
OOD nunca participan en la optimización de los modelos. Para cada conjunto de datos
se definen tres particiones no solapadas: \texttt{train} (muestras ID para el ajuste
de parámetros), \texttt{id\_eval} (muestras ID retenidas para el cálculo de métricas
y la construcción de la referencia de normalidad) y \texttt{ood\_eval} (muestras OOD
empleadas exclusivamente en la evaluación).

En los conjuntos sintéticos, el 70~\% de los puntos ID se asigna al entrenamiento
y el 30~\% restante a \texttt{id\_eval}, mientras que todos los puntos OOD forman el conjunto \texttt{ood\_eval}. Para MNIST y PathMNIST se emplean las particiones
\textit{train}/\textit{test} originales de los respectivos conjuntos. Para SICAP se
utilizan las divisiones proporcionadas por el departamento.

\subsection{Normalización}

Los conjuntos sintéticos y las características UNI de SICAP se normalizan con
estandarización~$z$-\textit{score} calculada por coordenada sobre el conjunto de
entrenamiento ($\tilde{x}_j = (x_j - \mu_j)/\sigma_j$, con denominadores
inferiores a $10^{-6}$ reemplazados por~$1$). Para MNIST se aplica la transformación
fija \texttt{Normalize}$(0.5,\,0.5)$ sobre la imagen en $[0,1]$, obteniendo el
rango~$[-1, 1]$. Las imágenes RGB de PathMNIST, almacenadas como enteros en
$[0, 255]$, se escalan directamente mediante $\tilde{x} = x/127.5 - 1$.

\subsection{Representación con modelo fundacional para SICAP}
\label{subsec:uni}

Las imágenes de SICAP no se procesan directamente como píxeles, sino que se codifican
primero con UNI. Este modelo fundacional produce un vector
de $1\,024$ dimensiones por parche que condensa información morfológica de alto nivel:
textura nuclear, densidad celular, arquitectura glandular y otros rasgos histológicos que
las redes entrenadas de cero con pocos datos difícilmente aprenden. Los vectores se
cargan de los ficheros precomputados \texttt{X\_train.npy} y \texttt{X\_test.npy}, de
modo que los modelos generativos que aprenden la distribución de SICAP operan íntegramente en el espacio de
características UNI y no en el de píxeles.

Esta decisión está motivada por razones prácticas y científicas. Desde el punto de
vista práctico, entrenar modelos generativos de calidad directamente sobre imágenes
histopatológicas de alta resolución con los recursos computacionales disponibles no
es viable. Desde el punto de vista científico, el espacio de características UNI
ofrece una representación compacta que facilita la comparación directa entre métodos
basados en modelos generativos y métodos de distancia directa en el espacio de
características (k-NN y Mahalanobis), dado que todos operan sobre la misma
representación de entrada.

% ─────────────────────────────────────────────────────────────────────────────
\section{Detalles de implementación}
\label{sec:exp-distancia}
% ─────────────────────────────────────────────────────────────────────────────

\subsection{Arquitecturas}
\label{subsec:arquitecturas}

Se implementan cuatro variantes de modelos generativos agrupadas en dos familias
(VAE y DDPM) y dos clases de \textit{backbone} (MLP residual y U-Net convolucional),
dando lugar a las combinaciones \textbf{MLP-VAE}, \textbf{MLP-DDPM}, \textbf{UNet-VAE}
y \textbf{UNet-DDPM}. Las variantes MLP se aplican a todos los conjuntos de datos;
las variantes UNet se reservan para los conjuntos de imágenes (MNIST y PathMNIST),
donde la estructura espacial de los datos justifica el uso de capas convolucionales.

\paragraph{MLP-VAE.}
El \emph{encoder} está formado por $d$ capas totalmente conectadas de anchura~$h$ con
activación LeakyReLU~($\alpha{=}0.2$), seguidas de dos proyecciones lineales paralelas
que producen la media $\boldsymbol{\mu}$ y el logaritmo de la varianza $\log
\boldsymbol{\sigma}^2$ del espacio latente de dimensión~$L$. El logaritmo de la varianza
se recorta al intervalo $[-10, 10]$ para evitar inestabilidades numéricas durante el
truco de la reparametrización. El \emph{decoder} es simétrico: $d$ capas de anchura~$h$
con activación ReLU seguidas de una proyección lineal a la dimensión de entrada. La
función de pérdida de reconstrucción es el error cuadrático medio (MSE), y el
término KL de la evidencia inferior (ELBO) se escala mediante un multiplicador~$\beta$
configurable.

\paragraph{MLP-DDPM.}
El \textit{denoiser} es un MLP residual que recibe la muestra
ruidosa $\mathbf{x}_t$ junto con el instante de tiempo $t$ codificado mediante
un \textit{embedding} sinusoidal de dimensión~$d_t$. Para conjuntos de baja dimensión
($\dim \leq 4$), se aplican adicionalmente \textit{embeddings} sinusoidales con
escala $25.0$ a cada coordenada de entrada por separado, técnica que mejora la
capacidad del modelo para capturar variaciones de alta frecuencia en espacios muy
reducidos~\cite{tancik_fourier_2020}. Cada bloque residual incorpora normalización
de capa (\texttt{LayerNorm}) antes de la proyección lineal y \textit{dropout}
configurable; la inicialización de pesos sigue el esquema Xavier uniforme. El
proceso de difusión sigue el \textit{schedule} \texttt{squaredcos\_cap\_v2} de la
biblioteca \texttt{diffusers}~\cite{von-platen-etal-2022-diffusers}, con valores de
$\beta \in [10^{-4}, 0.02]$, $T{=}1\,000$ pasos de entrenamiento y modo de predicción
del ruido ($\varepsilon$-predicción).

\paragraph{UNet-VAE.}
El UNet-VAE utiliza el módulo \texttt{AutoencoderKL} de la biblioteca
MONAI~\cite{cardoso_monai_2022}, un autoencoder variacional convolucional con
bloques residuales que opera directamente sobre imágenes $28\times28$. La arquitectura
dispone de tres niveles de resolución con canales $[32, 64, 128]$ y un bloque residual
por nivel; la atención no lineal está desactivada en todos los niveles puesto que las
resoluciones son demasiado pequeñas para que sea beneficiosa. El espacio latente tiene
$4$ canales y resolución espacial $7\times7$ (dos submuestreos de factor~2), lo que
produce un cuello de botella de $196$ dimensiones. La función de pérdida de
reconstrucción es la norma $L_1$, que produce reconstrucciones más nítidas que el MSE
para datos de imagen~\cite{zhao_loss_2017}.

\paragraph{UNet-DDPM.}
El UNet-DDPM utiliza el módulo \texttt{DiffusionModelUNet} de MONAI, un U-Net con
bloques residuales y mecanismo de atención de cabezas múltiples en los dos niveles de
mayor profundidad (canales $[64, 128, 256]$, \texttt{attention\_levels}~$= [\text{False},
\text{True}, \text{True}]$, cabezas de dimensión $[0, 64, 128]$). El \textit{scheduler}
emplea un \textit{schedule} lineal de betas (\texttt{linear\_beta}) de MONAI.

\subsection{Protocolo de entrenamiento}
\label{subsec:entrenamiento}

Todos los modelos se optimizan con \textbf{AdamW}~\cite{loshchilov_decoupled_2019}
para los conjuntos de mayor complejidad (SICAP, PathMNIST, MNIST) y con
\textbf{Adam}~\cite{kingma_adam_2015} para los conjuntos sintéticos. El decaimiento
de pesos (\textit{weight decay}) es $10^{-5}$ cuando se aplica AdamW. El
programador de la tasa de aprendizaje combina un \textbf{calentamiento lineal} durante
\texttt{warmup\_epochs} épocas (del $10\,\%$ al valor nominal) con un decaimiento
coseno hasta $\eta_{\min}{=}10^{-6}$; los gradientes se recortan a norma máxima~$1.0$
para estabilizar el entrenamiento de los modelos DDPM.

Para los modelos DDPM se activa la \textbf{media exponencial móvil} (\textbf{EMA})
de los parámetros con decaimiento~$0.999$: en cada paso, los pesos de la copia EMA
se actualizan como $\hat{\theta} \leftarrow 0.999\,\hat{\theta} + 0.001\,\theta$.
Al finalizar el entrenamiento, los parámetros del modelo se reemplazan por la copia
EMA. Esta estrategia suaviza las predicciones del \textit{denoiser} frente a los
últimos pasos del descenso de gradiente y mejora la consistencia de los \textit{scores}
OOD al reducir la varianza de las reconstrucciones.

El entrenamiento de los VAE incorpora un \textbf{calentamiento del peso KL}: el
multiplicador~$\beta$ del término KL de la ELBO crece linealmente desde~$0$ hasta su
valor objetivo a lo largo de las primeras \texttt{kl\_warmup\_epochs} épocas. Esta
estrategia previene el colapso del espacio latente durante las primeras etapas del
entrenamiento~\cite{bowman_generating_2016}, un fenómeno por el que el codificador
aprende a ignorar la entrada y el espacio latente colapsa hacia la distribución a priori,
haciendo inútiles los \textit{scores} latentes de OOD. El efecto es visible en las
figuras~\ref{fig:training_curves_vae_sicapc1} y~\ref{fig:training_curves_vae_sicapc12},
donde la KL sube de forma gradual y progresiva sin colapsar.

Todos los experimentos emplean \textit{early stopping} sobre la pérdida de validación
en \texttt{id\_eval}: si la pérdida no mejora durante \texttt{patience} épocas
consecutivas, el entrenamiento se detiene y se restauran los mejores pesos observados.
La tabla~\ref{tab:hparams} resume los hiperparámetros principales por familia de datos
y tipo de modelo; los hiperparámetros de difusión (\,$\beta_{\mathrm{start}}{=}10^{-4}$,
$\beta_{\mathrm{end}}{=}0.02$, $T{=}1\,000$) son comunes a todos los modelos DDPM.
El seguimiento de los experimentos se realiza con \textbf{Weights \&
Biases} bajo el proyecto \texttt{tfg-ood-diffusion}.

\begin{table}[htbp]
    \centering
    \caption{Hiperparámetros de los modelos por familia de datos. $h$~=~anchura
             de capa oculta (MLP), ch.~=~progresión de canales del U-Net
             (primer$\to$último nivel), $d$~=~profundidad, lat.~$L$~=~dimensión
             latente (VAE; UNet-VAE: 4~canales, resolución $7{\times}7$,
             196~dim.\ total), emb.~$d_t$~=~dimensión del \textit{embedding}
             temporal (MLP-DDPM), ep.~=~épocas máximas, bs.~=~\textit{batch size}.}
    \label{tab:hparams}
    \begin{tabular}{llrrrrrrr}
        \toprule
        Conjunto & Modelo & $h$ & ch. & $d$ & lat.~$L$ & emb.~$d_t$ & ep. & bs. \\
        \midrule
        Sintéticos        & MLP-VAE   & 32    & —              & 1 & 2    & —   & 80  & 256 \\
        Sintéticos        & MLP-DDPM  & 64    & —              & 2 & —    & 64  & 100 & 16  \\
        \midrule
        MNIST             & MLP-VAE   & 1024  & —              & 5 & 256  & —   & 200 & 256 \\
        MNIST             & MLP-DDPM  & 2048  & —              & 5 & —    & 512 & 400 & 256 \\
        MNIST             & UNet-VAE  & —     & $32{\to}128$   & 1 & 4ch  & —   & 150 & 256 \\
        MNIST             & UNet-DDPM & —     & $64{\to}256$   & 1 & —    & —   & 300 & 256 \\
        \midrule
        PathMNIST-C1      & MLP-VAE   & 1024  & —              & 5 & 256  & —   & 200 & 256 \\
        PathMNIST-C1      & MLP-DDPM  & 2048  & —              & 6 & —    & 512 & 500 & 128 \\
        PathMNIST-C1      & UNet-VAE  & —     & $32{\to}128$   & 1 & 4ch  & —   & 150 & 256 \\
        PathMNIST-C1      & UNet-DDPM & —     & $64{\to}256$   & 1 & —    & —   & 300 & 256 \\
        PathMNIST-C2      & UNet-VAE  & —     & $32{\to}128$   & 1 & 4ch  & —   & 150 & 256 \\
        PathMNIST-C2      & UNet-DDPM & —     & $64{\to}256$   & 1 & —    & —   & 300 & 256 \\
        \midrule
        SICAP~(C1/C12)    & MLP-VAE   & 512   & —              & 4 & 128  & —   & 150 & 256 \\
        SICAP~(C1/C12)    & MLP-DDPM  & 1024  & —              & 5 & —    & 256 & 700 & 512 \\
        \bottomrule
    \end{tabular}
\end{table}

\begin{figure}[htbp]
    \centering
    \includegraphics[width=\textwidth]{results-training-mlp_vae_sicap_c1_s42_lr0.0001-plots-training_curves.pdf}
    \caption{Evolución de la pérdida total, el error de reconstrucción~(MSE)
             y la divergencia KL durante el entrenamiento del VAE en SICAP-C1
             (ID: Gleason~1). El calentamiento gradual del peso~$\beta$
             (\texttt{kl\_warmup\_epochs}) se aprecia en la subida progresiva
             de la KL sin colapso del espacio latente.}
    \label{fig:training_curves_vae_sicapc1}
\end{figure}

\begin{figure}[htbp]
    \centering
    \includegraphics[width=\textwidth]{results-training-mlp_vae_sicap_c12_s42_lr0.0005-plots-training_curves.pdf}
    \caption{Curvas de entrenamiento del VAE en SICAP-C12 (ID: Gleason~1--2).
             La mayor cantidad de datos de entrenamiento respecto a SICAP-C1
             produce una convergencia más estable y menor varianza en la pérdida
             de reconstrucción.}
    \label{fig:training_curves_vae_sicapc12}
\end{figure}

\begin{figure}[htbp]
    \centering
    \includegraphics[width=\textwidth]{results-training-mlp_ddpm_sicap_c1_s42_lr0.0001-plots-training_curves.pdf}
    \caption{Evolución del MSE de predicción de ruido durante el entrenamiento
             del DDPM en SICAP-C1. El modelo emplea EMA y \textit{early stopping}
             con \texttt{patience}~$= 50$.}
    \label{fig:training_curves_ddpm_sicapc1}
\end{figure}

\begin{figure}[htbp]
    \centering
    \includegraphics[width=\textwidth]{results-training-mlp_ddpm_sicap_c12_s42_lr0.001-plots-training_curves.pdf}
    \caption{Curvas de entrenamiento del DDPM en SICAP-C12: MSE de predicción
             de ruido a lo largo de las épocas, con EMA y \textit{early stopping}
             activos.}
    \label{fig:training_curves_ddpm_sicapc12}
\end{figure}

% ─────────────────────────────────────────────────────────────────────────────
\section{Métricas de evaluación}
\label{sec:exp-metricas}
% ─────────────────────────────────────────────────────────────────────────────

La evaluación de la capacidad de detección OOD se plantea como un problema de
\textit{ranking} binario: dado un \textit{score}~$S(\mathbf{x})$, las muestras ID deben
recibir sistemáticamente valores más bajos que las muestras OOD. Se emplean tres
métricas agnósticas al umbral, ampliamente consolidadas en la literatura de detección
de distribuciones fuera de soporte~\cite{hendrycks2017baseline, yang_generalized_2024}:

\paragraph{AUROC (\textit{Area Under the Receiver Operating Characteristic}).}~\\
El AUROC mide la probabilidad de que una muestra OOD aleatoria reciba un
\textit{score} superior al de una muestra ID aleatoria. Su rango es $[0, 1]$: un
valor de~$1.0$ denota separación perfecta y~$0.5$ corresponde a un discriminador
aleatorio. Su invariancia al desbalanceo entre clases y su interpretación
probabilística directa la convierten en la métrica principal de este trabajo, con
la que se comparan y ordenan los distintos \textit{scores} y configuraciones.

\paragraph{AUPR (\textit{Area Under the Precision-Recall Curve}).}~\\
El AUPR resume la curva precisión--\textit{recall} bajo la convención de que la
clase positiva son las muestras OOD. A diferencia del AUROC, es sensible al
desbalanceo entre clases: decrece cuando el conjunto OOD es minoritario, lo que
la hace especialmente informativa en escenarios médicos donde el número de casos
patológicos puede ser mucho menor que el de casos normales. Se reporta como
métrica complementaria.

\paragraph{FPR@95TPR (\textit{False Positive Rate at 95\,\% TPR}).}~\\
El FPR@95TPR mide la proporción de muestras ID clasificadas erróneamente como OOD
cuando el umbral de decisión~$\tau$ se fija de modo que el $95\,\%$ de las muestras
OOD son detectadas correctamente. Se calcula localizando en la curva ROC el punto
de cruce con $\mathrm{TPR}{=}0.95$ y leyendo la FPR correspondiente. Valores bajos
indican que el modelo puede mantener una alta tasa de detección con pocos falsos
positivos sobre datos normales, lo que es especialmente relevante en aplicaciones
clínicas con alto coste de alarmas falsas.

\paragraph{Análisis de umbrales.}~\\
Las métricas anteriores se complementan con un análisis de umbral operativo: dado un
objetivo de FPR del $5\,\%$, el umbral~$\tau$ se calcula como el percentil~$95$ de
los \textit{scores} sobre \texttt{id\_eval}. Se reportan la TPR y FPR efectivos al
aplicar~$\tau$, lo que permite estimar el comportamiento del sistema en condiciones
de mundo abierto, donde un umbral fijo debe balancear la tasa de detección con la
tasa de falsas alarmas.

\endinput